% narrowband_weil.tex -- round 461, narrow-band Weil consistency probe.
% FORK_SHARPENED.
% Companion: verify_narrowband_weil.py.
% Discovery: experiments/tfpt-discovery/narrowband_weil_probe.py.
% NO RH CLAIM.  NO anti-RH claim.
\documentclass[11pt]{article}

\usepackage[a4paper,margin=2.55cm]{geometry}
\usepackage{amsmath,amssymb,amsthm}
\usepackage{microtype}
\usepackage{booktabs}
\usepackage{xcolor}
\usepackage[colorlinks=true,linkcolor=blue!50!black,urlcolor=blue!50!black]{hyperref}

\newtheorem{prop}{Proposition}
\newtheorem{remark}{Remark}
\newcommand{\QW}{Q_{\mathrm W}}
\newcommand{\Rd}{R^{\dagger}}

\title{\bfseries Narrow-band Weil audit\\[2mm]
\large What the frequently-selected mincut would prove}
\author{Stefan Hamann}
\date{August 30, 2026}

\begin{document}
\maketitle

\begin{abstract}\noindent
The frequently-selected Lean route does not currently prove RH.
Its formal endpoint is nonnegativity of a three-channel object
\texttt{weilForm} on dyadic piecewise-linear autocorrelations.
The spectral/zero side and the density/continuity passage to the
standard Weil criterion are explicitly absent; the path also
consumes one typed classical \texttt{sorryAx} and one named,
unproved Hankel/read bridge.  After those gaps are closed, the
correct classical translation is not a harmless narrow-band
statement: it is positivity of the compact-window Weil form on
an unbounded cofinal sequence of support windows, hence---by
nestedness---the full Weil criterion.

Numerically, the r459 selected builder has effective
autocorrelation support
\[
 U_k=N_k\Delta_k=\frac{k\log2}{2}.
\]
Every measured member $k=5,\ldots,12$ lies beyond both the
classical prime-blind range $U<\log2$ and the recent
arXiv:2608.24827 certificate $U\le1.6$.  The r459 race
$0.61$--$0.76$ therefore measures a genuinely conditional tail,
subject to the still-open dictionary.  Verdict:
\texttt{FORK\_SHARPENED}.  No RH claim.
\end{abstract}

\medskip\noindent
\textbf{Status.}\enspace
Consistency and solution probe only.
No new theorem about $\zeta$.
No $L^*$ claim and no $\Rd$ claim.
No RH claim and no anti-RH claim.

\section{Q1: the statement actually present in Lean}

\subsection{The test class}

\begin{prop}[Exact Lean class]\label{prop:class}
A \texttt{GridElement} consists of
\[
 (n_g,m_f,x),\qquad x:\{0,\ldots,n_g-1\}\longrightarrow\mathbb R,
 \qquad D_0=2^{-m_f}.
\]
It defines the autocorrelation coefficients
\[
 a_d=D_0\sum_i x_i x_{i+d}
\]
and the even piecewise-linear interpolant $g_x$ through
$a_d$ on the dyadic grid.  Its support is contained in
$[-n_gD_0,n_gD_0]$.
\end{prop}

\noindent
This is literal from
\texttt{Elementwise.lean}:248--253 (fields),
260--273 ($D_0$ and autocorrelation),
293--301 (even piecewise-linear interpolant), and
303--334 (support and evenness).
There is no smoothness field.  The induced $g_x$ is continuous
and piecewise linear, generally not $C^\infty$.

Classically $g_x=h_x*\widetilde h_x$, where $h_x$ is the
corresponding dyadic step function.  Therefore
\[
 \widehat g_x(t)=|\widehat h_x(t)|^2\ge0.
\]
Thus the relevant positivity is Fourier positivity of the
\emph{autocorrelation}; it is not an additional assumption on
an arbitrary test function.

\subsection{The three-channel value}

The arithmetic channel is genuinely transcribed:
\[
\begin{split}
\texttt{combRead}(a,g_x)
  &=\sum_{n\in\texttt{windowAtoms}(a)}
       \frac{2\Lambda(n)}{\sqrt n}\,g_x(\log n),\\
\texttt{fullRead}(a,m,g_x)
  &=\texttt{archRead}-\texttt{combRead}+\texttt{poleRead}.
\end{split}
\]
See \texttt{Elementwise.lean}:361--391 and 711--728.
The exact finite stabilization of the comb is proved at
\texttt{Elementwise.lean}:396--430.
The archimedean read and its Weil pairing remain opaque
(\texttt{Elementwise.lean}:553--565), and their identification is
the typed \texttt{sorry}
\path{arch_elementwise_stabilization}
(\texttt{Elementwise.lean}:676--698).

\subsection{Selected covering and endpoint}

The sequence is exactly
\[
 a_k=2^k,\qquad r_k=\lfloor\sqrt k\rfloor,\qquad
 m_k=k2^{r_k}-1,\qquad
 \Delta_k=2^{-r_k}\log2,
\]
at \texttt{Selected.lean}:293--343.
Both $a_k$ and $m_k$ tend to infinity and
\texttt{selected\_covers} supplies, for each fixed grid element,
one sufficiently large selected index
(\texttt{Selected.lean}:347--418).
There is no fixed support ceiling in this quantifier: the
\texttt{steps} field is unbounded and the sequence is cofinal.

The extraction
\path{weil_nonneg_of_frequently_plain}
(\texttt{FrequentlySelected.lean}:345--355) takes one index that
is both covering and good.  The named mincut is
\[
 \exists^\infty k:\quad
 \Rd(W_k^{\mathbb R})\succeq\tfrac12 I
\]
(\texttt{FrequentlySelected.lean}:260--285).
The real-window Loewner equivalence to
$A_{\mathrm{cap}}\succeq0$ is proved
(\texttt{FrequentlySelected.lean}:370--380).
The remaining implication
\[
 A_{\mathrm{cap}}\succeq0
 \Longrightarrow
 \texttt{fullRead}(a_k,m_k,g_x)\ge0
\]
is only the named Prop
\[
 \text{\path{SelectedACapPsdImpliesPlainReads}}
\]
at \texttt{FrequentlySelected.lean}:382--409.

\begin{prop}[Formal endpoint]\label{prop:endpoint}
Even after assuming the mincut and the named bridge, the theorem
historically called \path{rh_of_frequently_selected} concludes
only
\[
 \forall f:\texttt{GridElement},\qquad
 0\le\texttt{weilForm}(f).
\]
It does not conclude the Riemann hypothesis.
\end{prop}

\noindent
See \texttt{FrequentlySelected.lean}:430--449.
The source explicitly says that the spectral/zero side is not
formalized (\texttt{Elementwise.lean}:61--77, 721--726), and that
the density/continuity step off the native class is deliberately
not formalized (\texttt{Elementwise.lean}:88--93).
The audit records \texttt{sorryAx} on this path and names the
read-identification remainder
(\texttt{Audit.lean}:280--293, 323--339).
Consequently the premise ``Lean mincut $\Rightarrow$ RH,
sorry-free'' is false in the current tree.

\subsection{A second dictionary warning}

The operational r459 builder sets
$M=m_k+1$ and $N_k=M/2$.
The Lean object \texttt{ExactFold}, however, first maps exact
$\log n$ nodes by a cosine and merges equal images
(\texttt{Selected.lean}:160--179;
\texttt{Source.lean}:400--424); its \texttt{cap} is
$(S+1)/2$ for that exact folded support
(\texttt{Selected.lean}:242--246).
The equality of this exact-fold cap with the r459
$M/2$ tent/DFT cap is not a theorem.
This is part of the same open Hankel/read dictionary and must be
closed before builder numerics can be promoted to the Lean mincut.

\section{Q2: the unconditionedness fork}

\subsection{What is unconditionally known}

Use the additive-log form of Weil's criterion.  For a compactly
supported base function $h$, let
$g=h*\widetilde h$.  Weil positivity is
\[
 \QW(h):=W(g)\ge0.
\]
For $h\in C_c^\infty(\mathbb R)$, positivity for every $h$ is
equivalent to RH (Weil 1952; modern formulations use this
nonnegative-definiteness of the Weil distribution).

Bombieri's Theorem~12 proves strict positivity when the support
of $h$ is contained in an interval of length
\[
 U<\log2.
\]
After translation take
$\operatorname{supp}h\subset[-U/2,U/2]$; then
$\operatorname{supp}(h*\widetilde h)\subset[-U,U]$, so the
von-Mangoldt sum vanishes.  This is the precise
Yoshida/Bombieri small-support result, not a Fourier-support
statement with the number $2$.

A recent preprint, arXiv:2608.24827 (August 25, 2026), claims the
certified bound
\[
 \QW(h)\ge8.9\cdot10^{-18}\|h\|_2^2,\qquad
 \operatorname{supp}h\subset[-0.8,0.8],
\]
hence autocorrelation support $U\le1.6$.
This extends the classical support length $\log2$ but is a
recent computational preprint, not the classical
Yoshida/Bombieri theorem.

\begin{remark}[The alleged $(-2,2)$ theorem is the wrong regime]
There is no standard unconditional theorem saying that
``$\operatorname{supp}\widehat f\subset(-2,2)$ implies Weil
positivity for $\zeta$'' in the normalization above.
Support restrictions such as $(-2,2)$ occur in scaled
one-/two-level density and pair-correlation formulae; their
coordinate is normalized by a conductor or height.
They are not a universal compact-log-support certificate for
the Weil quadratic form.  Fej\'er kernels supply
nonnegative Fourier transforms, but do not make the prime term
of the Weil formula nonnegative for arbitrary support.
\end{remark}

\subsection{Numerical location of the selected windows}

In the r459 reference builder,
\[
 N_k=\frac{m_k+1}{2}=k2^{r_k-1},
\qquad U_k=N_k\Delta_k=\frac{k\log2}{2}.
\]
$U_k$ is the support half-width of the autocorrelation, equivalently
the length of the underlying step-function interval.
In the centered convention
$\operatorname{supp}h\subset[-L_k,L_k]$,
\[
 L_k=\frac{U_k}{2}=\frac{k\log2}{4}.
\]

\begin{center}
\small
\begin{tabular}{@{}rrrrrrrrrr@{}}
\toprule
$k$ & $m_k$ & $N_k$ & $\Delta_k$ & $U_k$ & $L_k$ &
$J_P$ & $n_{1.6}$ & tail & race\\
\midrule
5  & 19 & 10 & .1732868 & 1.732868 & .866434 & 3 & 9  & 1  & .675060\\
6  & 23 & 12 & .1732868 & 2.079442 & 1.039721 & 3 & 9  & 3  & .694907\\
7  & 27 & 14 & .1732868 & 2.426015 & 1.213008 & 3 & 9  & 5  & .628736\\
8  & 31 & 16 & .1732868 & 2.772589 & 1.386294 & 3 & 9  & 7  & .606566\\
9  & 71 & 36 & .0866434 & 3.119162 & 1.559581 & 7 & 18 & 18 & .635369\\
10 & 79 & 40 & .0866434 & 3.465736 & 1.732868 & 7 & 18 & 22 & .757783\\
11 & 87 & 44 & .0866434 & 3.812309 & 1.906155 & 7 & 18 & 26 & .715782\\
12 & 95 & 48 & .0866434 & 4.158883 & 2.079442 & 7 & 18 & 30 & .693564\\
\bottomrule
\end{tabular}
\end{center}

\noindent
Here $J_P=2^{r_k}-1$ is the r459 prime-blind prefix marker.
$n_{1.6}=\lfloor1.6/\Delta_k\rfloor$ is the last depth covered
by the recent $U\le1.6$ certificate, and ``tail'' is
$N_k-n_{1.6}$.
Every measured member is outside the recent range:
the crossing occurs between $k=4$ ($U_4=1.38629$) and
$k=5$ ($U_5=1.73287$).
The classical crossing occurs already between $k=2$
($U_2=\log2$ at the endpoint) and $k=3$.

\begin{prop}[Fork verdict]\label{prop:fork}
Alternative (a) does not occur on the measured ladder.
Alternative (b) starts at its first member, $k=5$.
The r459 race is therefore a finite measurement in the
conditional support tail, not a consequence of the classical
narrow-band theorem.
\end{prop}

\subsection{Literal classical translation of the mincut}

\begin{quote}
\textbf{Literal translation.}
For infinitely many unbounded selected indices $k$, the augmented
dual-resolvent cone of the full selected prime window is
semidefinite at the half-filling cap.  If the missing
Hankel/read identification is valid, this says that the
Riemann--Weil quadratic functional is nonnegative on every
dyadic step-function autocorrelation resolved and covered by
that selected window.  Since the selected supports are cofinal,
every compactly supported dyadic autocorrelation is covered by
one good window.  After the missing continuity/density and
spectral explicit-formula identifications, this is
\emph{cofinal compact-window Weil positivity}.
\end{quote}

Define the classical compact-window floor
\[
 \lambda_*(L)=
 \inf_{\substack{0\ne h\in L^2(\mathbb R)\\
                  \operatorname{supp}h\subset[-L,L]}}
 \frac{\QW(h)}{\|h\|_2^2}.
\]
The honest classical name for the completed mincut is
\[
 \lambda_*(L_k)\ge0\quad\hbox{for an unbounded sequence }L_k.
\]
Because the spaces are nested, $\lambda_*(L)$ is nonincreasing.
Hence positivity on any unbounded cofinal subsequence implies
positivity on every finite $L$.  This is not a weaker
``positive-density'' fragment after the dictionary is closed;
it is the full Weil criterion.

\section{Q3: best complete route}

\begin{center}
\small
\begin{tabular}{@{}p{.06\textwidth}p{.59\textwidth}p{.25\textwidth}@{}}
\toprule
Step & Task & Honest status\\
\midrule
1 &
Prove that the exact folded Lean window and the tent/DFT r459
window have the same relevant Hankel form and cap; repair any
dimension or mesh-normalization mismatch. &
\textbf{Open, local and decisive}.  Finite algebra after a faithful
fold transcription, but not present.\\
\addlinespace
2 &
Prove \texttt{SelectedACapPsdImpliesPlainReads}, including
bordered-to-plain compression and the exact arch/border read. &
\textbf{Open, plausibly mechanizable}; this is the current
load-bearing bridge, not the $\Rd$ Loewner algebra.\\
\addlinespace
3 &
Replace \texttt{arch\_elementwise\_stabilization}'s
\texttt{sorry} by the classical digamma/Mellin identity. &
\textbf{Known mathematics, formalization work}.\\
\addlinespace
4 &
Formalize dyadic step density, continuity of the localized Weil
form, and the Guinand--Weil spectral/zero identity. &
\textbf{Known mathematics, substantial Lean work}.  Only after
this step may the endpoint be called an RH criterion.\\
\addlinespace
5 &
Certify $\lambda_*(L)\ge0$ at additional fixed $L$ by
interval matrices.  Reproduce the claimed $L=.8$ certificate,
then target $L_5=.866434$. &
\textbf{Finite and in principle feasible}; margins may be extremely
small.  Each fixed $L$ is only a finite RH fragment.\\
\addlinespace
6 &
Prove nonnegativity on unbounded $L_k$ (equivalently, the
frequently-selected mincut after Steps 1--4). &
\textbf{Open-hard: equivalent to RH}.  No known density or
subsequence theorem supplies it.\\
\bottomrule
\end{tabular}
\end{center}

\paragraph{Why zero-density theorems do not close Step 6.}
Results that put a positive proportion of zeros on the critical
line control averages or restricted zero sets.  Weil positivity
requires nonnegativity against every compactly supported
autocorrelation.  A single off-line zero can create a negative
direction once the support space is large enough.  Thus
``many zeros on the line'' does not imply the needed cofinal
window positivity.

\paragraph{Why Li/Bombieri--Lagarias do not close Step 6.}
Li's theorem says RH is equivalent to
$\lambda_n^{\mathrm{Li}}>0$ for every $n\ge1$.
Bombieri--Lagarias derive these coefficients from the
Guinand--Weil functional and generalize the zero-set criterion.
This is an exact coordinate bridge, useful for cross-checking
normalizations and constructing test vectors.  It does not make
the positivity unconditional; proving all Li coefficients
positive is again RH.

\paragraph{Useful diagnostic result.}
Bombieri's variational analysis shows that, under a finite
off-line-zero hypothesis, sufficiently rich truncations acquire
negative eigenvalues whose count detects the failures.
This is valuable for red-teaming the selected matrices, but it is
not a positive lower bound on a cofinal family.

\section{Conclusion}

The chain is consistent because it is not presently a closed RH
proof and because its measured windows are not in the
unconditional narrow-band regime.  The strongest honest
externalization is:
\begin{quote}
The selected $\Rd$ mincut is a proposed finite-matrix realization
of cofinal compact-window Weil positivity.  Its Loewner algebra is
formalized, but the exact Hankel/read dictionary, one classical
archimedean identity, and the density/spectral completion remain
open.  On the r459 builder the tested windows begin just beyond
the current $U=1.6$ compact-support certificate and stay positive
with race $0.61$--$0.76$ through $k=12$.
\end{quote}

\noindent
\textbf{Verdict: \texttt{FORK\_SHARPENED}.}
No RH claim.  No anti-RH claim.

\begin{thebibliography}{9}
\bibitem{Weil1952}
A.~Weil, \emph{Sur les formules explicites de la th\'eorie des
nombres premiers}, Meddelanden Fr\r an Lunds Univ. Mat. Sem.,
1952, 252--265.

\bibitem{Yoshida1992}
H.~Yoshida, \emph{On Hermitian forms attached to zeta functions},
Adv. Stud. Pure Math. \textbf{21} (1992), 281--325.

\bibitem{Bombieri2000}
E.~Bombieri, \emph{Remarks on Weil's quadratic functional in the
theory of prime numbers, I}, Rend. Lincei Mat. Appl.
\textbf{11} (2000), 183--233; Theorem~12.

\bibitem{ConnesConsani2021}
A.~Connes and C.~Consani, \emph{Weil positivity and trace formula,
the archimedean place}, Selecta Math. \textbf{27} (2021), Paper 77.

\bibitem{Li1997}
X.-J.~Li, \emph{The positivity of a sequence of numbers and the
Riemann hypothesis}, J. Number Theory \textbf{65} (1997), 325--333.

\bibitem{BombieriLagarias1999}
E.~Bombieri and J.~C.~Lagarias,
\emph{Complements to Li's criterion for the Riemann hypothesis},
J. Number Theory \textbf{77} (1999), 274--287.

\bibitem{Chuk2026}
M.~Chuk, \emph{Weil positivity in compact windows: certified
two-sided bounds and a Landau--Widom decay law},
arXiv:2608.24827 (August 25, 2026), preprint.
\end{thebibliography}

\end{document}
